\begin{ex}
Số 4725 gồm mấy nghìn?
\shortans{4}
\loigiai{
Chữ số hàng nghìn là 4.
}
\end{ex}

\begin{ex}
Số liền sau của 6999 là số nào?
\shortans{7000}
\loigiai{
$6999 + 1 = 7000$.
}
\end{ex}

\begin{ex}
Số liền trước của 8000 là số nào?
\shortans{7999}
\loigiai{
$8000 - 1 = 7999$.
}
\end{ex}

\begin{ex}
Trong các số $6450, 6540, 6405, 6504$, số lớn nhất là số nào?
\shortans{6540}
\loigiai{
Số lớn nhất là 6540.
}
\end{ex}

\begin{ex}
Số La Mã XVIII tương ứng với số tự nhiên nào?
\shortans{18}
\loigiai{
XVIII là 18.
}
\end{ex}

\begin{ex}
Số La Mã XIV tương ứng với số tự nhiên nào?
\shortans{14}
\loigiai{
XIV là 14.
}
\end{ex}

\begin{ex}
Làm tròn số 564 đến hàng chục được số nào?
\shortans{560}
\loigiai{
Làm tròn xuống 560.
}
\end{ex}

\begin{ex}
Làm tròn số 750 đến hàng trăm được số nào?
\shortans{800}
\loigiai{
Làm tròn lên 800.
}
\end{ex}

\begin{ex}
Viết số thích hợp: $7000 + 300 + 4 = \dots$
\shortans{7304}
\loigiai{
$7000 + 300 + 4 = 7304$.
}
\end{ex}

\begin{ex}
Tìm $x$ là số tròn nghìn biết $4000 < x < 6000$.
\shortans{5000}
\loigiai{
Số đó là 5000.
}
\end{ex}

\begin{ex}
Cho 4 chữ số 2, 0, 8, 9. Số lớn nhất có 4 chữ số ghép từ cả 4 chữ số này là số nào?
\shortans{9820}
\loigiai{
Số đó là 9820.
}
\end{ex}

\begin{ex}
Cho 4 chữ số 1, 0, 5, 8. Số bé nhất có 4 chữ số khác nhau ghép từ 4 chữ số này là số nào?
\shortans{1058}
\loigiai{
Số đó là 1058.
}
\end{ex}

\begin{ex}
Khoảng cách giữa hai số 9000 và 8700 là bao nhiêu?
\shortans{300}
\loigiai{
$9000 - 8700 = 300$.
}
\end{ex}

\begin{ex}
Đồng hồ La Mã chỉ vào số VIII nghĩa là mấy giờ?
\shortans{8}
\loigiai{
VIII là 8 giờ.
}
\end{ex}

\begin{ex}
Số 5555 có bao nhiêu chữ số 5?
\shortans{4}
\loigiai{
Có 4 chữ số 5.
}
\end{ex}

\begin{ex}
Làm tròn số 8263 đến hàng trăm ta được số nào?
\shortans{8300}
\loigiai{
Hàng chục là 6 nên làm tròn lên 8300.
}
\end{ex}

\begin{ex}
Làm tròn số 4096 đến hàng chục ta được số nào?
\shortans{4100}
\loigiai{
Hàng đơn vị là 6 làm tròn lên: $4090 + 10 = 4100$.
}
\end{ex}

\begin{ex}
Số tự nhiên nhỏ nhất có 4 chữ số khác nhau là số nào?
\shortans{1023}
\loigiai{
Số đó là 1023.
}
\end{ex}

\begin{ex}
Số tự nhiên lớn nhất có 4 chữ số khác nhau là số nào?
\shortans{9876}
\loigiai{
Số đó là 9876.
}
\end{ex}

\begin{ex}
Tính: XXI - IX = \dots (kết quả bằng số tự nhiên).
\shortans{12}
\loigiai{
$21 - 9 = 12$.
}
\end{ex}

\begin{ex}
Tính: VI $\times$ III = \dots (kết quả bằng số tự nhiên).
\shortans{18}
\loigiai{
$6 \times 3 = 18$.
}
\end{ex}

\begin{ex}
Số lớn nhất có 4 chữ số có tích các chữ số bằng 1 là số nào?
\shortans{1111}
\loigiai{
Tất cả các chữ số phải là 1: 1111.
}
\end{ex}

\begin{ex}
Có bao nhiêu số tròn trăm từ 2000 đến 3000?
\shortans{11}
\loigiai{
2000, 2100, ..., 3000 (11 số).
}
\end{ex}

\begin{ex}
Làm tròn số 649 đến hàng trăm được số nào?
\shortans{600}
\loigiai{
Hàng chục là 4 nên làm tròn thành 600.
}
\end{ex}

\begin{ex}
Làm tròn số 650 đến hàng trăm được số nào?
\shortans{700}
\loigiai{
Hàng chục là 5 nên làm tròn lên 700.
}
\end{ex}
